\documentclass[a4paper,11pt]{article}

\usepackage[spanish]{babel}
\usepackage[utf8]{inputenc}
\usepackage[T1]{fontenc}
\usepackage{geometry}
\usepackage{tikz}
\usetikzlibrary{automata, positioning, arrows.meta, fit, backgrounds, shapes.geometric}
\usepackage{longtable}
\usepackage{array}
\usepackage{booktabs}
\usepackage{colortbl}
\usepackage{xcolor}
\usepackage{enumitem}
\usepackage{tabularx}

\geometry{margin=2cm}

\begin{document}

% ============================================================
\section*{1.\; Arquitectura General del Sistema-on-Chip (SoC)}
% ============================================================

El dise\~no del subsistema se rige por un esquema \textbf{Multi-Maestro (Multi-Master)} donde el procesador MIPS y el M\'odulo de Entrada/Salida (E/S) compiten por el control del Bus del Sistema recurriendo al \'Arbitro.

A continuaci\'on se presenta el diagrama de rutas de datos y peticiones del sistema:

\bigskip
\begin{center}
\shorthandoff{<>}
\resizebox{\linewidth}{!}{%
\begin{tikzpicture}[
    >=Stealth,
    shorten >=1pt,
    font=\normalsize,
    box/.style={
        draw, very thick, rectangle,
        minimum width=4.5cm, minimum height=1.3cm,
        align=center, fill=white, rounded corners=3pt
    },
    grpbox/.style={
        draw, very thick, rounded corners=8pt,
        inner sep=22pt, fill=gray!4
    },
    lbl/.style={font=\footnotesize, align=center, fill=white,
                inner sep=2pt, rounded corners=2pt}
]

%% ---- Coordenadas absolutas ----
%%  Fila A (Maestros):  y = 9
%%  Fila B (Control):   y = 4.5
%%  Fila C (Esclavos):  y = 0

\node[box] (ES)     at (0,   9)   {M\'odulo E/S};
\node[box] (MIPS)   at (13,  9)   {Procesador MIPS};

\node[box] (Arb)    at (0,   4.5) {\'Arbitro del Bus};
\node[box] (Cache)  at (6.5, 4.5) {Controlador\\de Cach\'e};
\node[box] (Decod)  at (13,  4.5) {Decodificador\\(IF Dir)};

\node[box] (RAM)    at (3.5, 0)   {Memoria Principal\\(RAM)};
\node[box] (Scratch)at (10.5,0)   {Memoria\\Scratchpad};

%% ---- Grupos ----
\begin{scope}[on background layer]
  \node[grpbox, fit=(ES)(MIPS),
        label={[font=\small\bfseries]above:Maestros / Productores}] {};
\end{scope}
\begin{scope}[on background layer]
  \node[grpbox, fit=(RAM)(Scratch),
        label={[font=\small\bfseries]below:Esclavos / Memorias}] {};
\end{scope}

%% ---- Flujos de control (todas lineas rectas) ----
% ES -> Arbitro (recto vertical)
\draw[->] (ES.south) -- node[lbl, left] {Petici\'on Bus\\para E/S} (Arb.north);

% MIPS -> Decod (recto vertical)
\draw[->] (MIPS.south) -- node[lbl, right] {Petici\'on\\(Dir, Dato)} (Decod.north);

% Decod -> Cache (recto horizontal)
\draw[->] (Decod.west) -- node[lbl, above] {Rango RAM} (Cache.east);

% Decod -> Scratch: baja vertical y luego horizontal a Scratch
\draw[->] (Decod.south) -- ++(0,-1.5) -|
      node[lbl, right, pos=0.25] {Rango Scratchpad}
      (Scratch.north);

% Cache -> Arbitro (recto horizontal)
\draw[->] (Cache.west) -- node[lbl, above] {Petici\'on Bus\\(Fallo/CopyBack)} (Arb.east);

%% ---- Retornos al MIPS (lineas rectas ortogonales) ----
% Cache -> MIPS: sube vertical hasta y=7.5, luego horizontal hasta MIPS
\draw[->] (Cache.north) -- ++(0,2.2) coordinate (hitWP)
      -- node[lbl, above] {Hit=1 (Acierto)}
      (hitWP -| MIPS.south west) -- (MIPS.south west);

% Scratch -> MIPS: sube vertical (misma columna x=10.5 que Scratch)
\draw[->] (Scratch.north) -- ++(0,1.2) coordinate (respWP)
      -| node[lbl, pos=0.5] {Respuesta Directa}
      (MIPS.south);

%% ---- Transferencias BUS (lineas rectas ortogonales) ----
% Cache <-> RAM: baja vertical y luego horizontal
\draw[<->, thick, dashed] (Cache.south) -- ++(0,-1.2) coordinate (busC)
      -| node[lbl, pos=0.5] {[BUS] Transferencia\\bloques}
      (RAM.north);

% ES <-> RAM: sale izquierda de ES (x=-2), baja vertical, entra por RAM.west
\draw[<->, thick, dashed]
      (ES.west) -- ++(-2, 0) coordinate (viaL)
      -- node[lbl, left] {[BUS]\\E/S$\leftrightarrow$RAM}
      (viaL |- RAM.west) -- (RAM.west);

% ES <-> Scratch: sube de ES (y=11), va horizontal derecha, baja a Scratch
\draw[<->, thick, dashed]
      (ES.north) -- ++(0, 1.5) coordinate (viaT)
      -- node[lbl, above] {[BUS] E/S$\leftrightarrow$Scratch}
      (viaT -| Scratch.north) -- (Scratch.north);

\end{tikzpicture}%
}
\shorthandon{<>}
\end{center}
\bigskip

% ============================================================
\section*{2.\; Pol\'iticas de Configuraci\'on de la Memoria Cach\'e}
% ============================================================

\begin{itemize}[leftmargin=1.5em]
  \item \textbf{Escritura en Aciertos (Write Hit): Copy-Back.}
        Los datos modificados se escriben \'unicamente en la V\'ia de la Cach\'e elevando el \textit{Dirty Bit} (`Sucio'). Esta pol\'itica optimiza el tr\'afico, ya que no se accede al bus exterior innecesariamente hasta que el bloque sea forzosamente desalojado.
  \item \textbf{Escritura en Fallos (Write Miss): Write-Around (Write-No-Allocate).}
        Si la direcci\'on de memoria no se encuentra en las v\'ias, el controlador rechaza su ingesta. Confiere una petici\'on de Bus directa a la RAM principal, y escribe la palabra sin alterar en ning\'un momento el contenido interno actual de la Cach\'e.
\end{itemize}

\bigskip

% ============================================================
\section*{3.\; Matriz Gen\'erica de Casos de Uso}
% ============================================================

Teniendo en cuenta las rutas del datapath y las pol\'iticas del MIPS y MC, la memoria lidia de la siguiente manera con los 9 escenarios est\'andar:

\renewcommand{\arraystretch}{1.25}

\begin{longtable}{|p{0.45cm}|p{4.0cm}|p{1.6cm}|p{2.2cm}|p{3.5cm}|p{1.0cm}|p{4.0cm}|}
\hline
\textbf{ID} &
\textbf{Operaci\'on / Escenario} &
\textbf{Maestro} &
\textbf{Destino} &
\textbf{Condici\'on} &
\textbf{Bus?} &
\textbf{Flujo resultante} \\
\hline
\endhead
\hline
\endfoot
1 & MIPS Acceso a Scratchpad     & MIPS    & MemScratch & Address en rango dedicado.                           & No  & Petici\'on directa por IF. La Scratchpad responde en 1 ciclo. \\
\hline
2 & Acierto LECTURA Cach\'e      & MIPS    & Cach\'e     & Instr. \texttt{lw} / Hit=1.                          & No  & El MC verifica Tag+Validez y responde en 1 ciclo. \\
\hline
3 & Acierto ESCRITURA Cach\'e    & MIPS    & Cach\'e     & Instr. \texttt{sw} / Hit=1.                          & No  & Se sobreescribe la palabra y se levanta Dirty Bit=1. \\
\hline
4 & Fallo de LECTURA Cach\'e     & Cach\'e  & RAM         & Instr. \texttt{lw} / Hit=0.                          & S\'i & MC detiene pipeline, pide bus, transfiere bloque a la v\'ia y presenta el dato. \\
\hline
5 & Fallo ESCRITURA Cach\'e      & Cach\'e  & RAM         & Instr. \texttt{sw} / Hit=0.                          & S\'i & \textit{Write-Around}: Sin alocar. Pide bus y escribe directamente en RAM. \\
\hline
6 & Desalojo de Bloque (CopyBack)& Cach\'e  & RAM         & Fallo que precisa v\'ia con DirtyBit=1.               & S\'i & Pide bus, vuelca bloque sucio a RAM y reanuda la petici\'on pendiente. \\
\hline
7 & DMA E/S a M. Principal       & M\'od.\ E/S & RAM      & Actividad de E/S externa.                             & S\'i & \'Arbitro cede bus al perif\'erico; transmite r\'afagas a RAM. \\
\hline
8 & DMA E/S a MemScratch         & M\'od.\ E/S & Scratchpad& Datos externos cr\'iticos (GPU/Red).                  & S\'i & \'Arbitro cede bus; datos DMA van a Scratchpad permitiendo al MIPS computar en L1. \\
\hline
9 & Colisi\'on de Maestros        & Cach\'e/E/S & \'Arbitro  & Petici\'on sincr\'onica simult\'anea.                  & S\'i & \'Arbitro detecta doble \texttt{Req}, aplica Round-Robin/Prioridad y concede \texttt{Grant} al ganador. \\
\hline
\end{longtable}

\bigskip

% ============================================================
\section*{UC de la cach\'e: diagrama y tabla de estados}

\subsection*{Diagrama de estados}

\begin{center}
\shorthandoff{<>}
\begin{tikzpicture}[
    >=Stealth,
    shorten >=1pt,
    node distance=4.4cm and 4.0cm,
    on grid,
    auto,
    font=\small,
    every state/.style={
        draw,
        thick,
        minimum width=3.4cm,
        minimum height=1.15cm,
        align=center
    }
]

\node[state, initial] (inicio) {Inicio};
\node[state, right=of inicio] (pedir) {PedirBus};
\node[state, right=of pedir] (enviar) {EnviarAddr};
\node[state, above=of enviar] (copyback) {CopyBack};
\node[state, right=of enviar] (fetch) {FetchBloque};
\node[state, below=of enviar] (transfer) {TransferPalabra};
\node[state, below=of fetch] (fin) {FinAcceso};
\node[state, below=of pedir] (errorbus) {ErrorBus};

\path[->]
(inicio) edge[loop above] node[font=\footnotesize] {sin acceso} (inicio)
(inicio) edge[loop below] node[font=\footnotesize] {hit, error local,\\ registro interno} (inicio)
(inicio) edge node[above, font=\footnotesize] {miss / no cacheable} (pedir)

(pedir) edge[loop above] node[font=\footnotesize] {$grant=0$} (pedir)
(pedir) edge node[above, font=\footnotesize] {$grant=1$} (enviar)

(enviar) edge node[left, font=\footnotesize] {$DevSel=0$} (errorbus)
(enviar) edge node[below, font=\footnotesize] {non-cacheable} (transfer)
(enviar) edge[bend left=20] node[above, font=\footnotesize] {cacheable y dirty} (copyback)
(enviar) edge node[above, font=\footnotesize] {cacheable y clean} (fetch)

(copyback) edge[loop above] node[font=\footnotesize] {$TRDY=0$ o no última} (copyback)
(copyback) edge[bend left=20] node[right, font=\footnotesize] {última palabra} (enviar)

(fetch) edge[loop above] node[font=\footnotesize] {$TRDY=0$ o no última} (fetch)
(fetch) edge node[right, font=\footnotesize] {última palabra} (fin)

(transfer) edge[loop below] node[font=\footnotesize] {$TRDY=0$} (transfer)
(transfer) edge node[right, font=\footnotesize] {TRDY=1} (fin)

(fin) edge[bend left=20] node[below, font=\footnotesize] {fin} (inicio)
(errorbus) edge[bend right=20] node[left, font=\footnotesize] {error} (inicio);

\end{tikzpicture}
\shorthandon{<>}
\end{center}

\subsection*{Resumen de estados}

\renewcommand{\arraystretch}{1.2}

\begin{longtable}{|p{2.2cm}|p{4.4cm}|p{2.5cm}|p{5.8cm}|}
\hline
\textbf{Estado} & \textbf{Condición} & \textbf{Siguiente} & \textbf{Señales clave} \\
\hline
Inicio & sin acceso & Inicio & $ready=1$ \\
\hline
Inicio & hit de lectura & Inicio & $ready=1$, $inc\_r=1$, $mux\_output=00$ \\
\hline
Inicio & hit de escritura & Inicio & $ready=1$, $MC\_WE0/1$, $Update\_dirty=1$, $inc\_w=1$ \\
\hline
Inicio & miss cacheable o acceso no cacheable & PedirBus & $Bus\_req=1$ \\
\hline
PedirBus & $grant=0$ & PedirBus & $Bus\_req=1$ \\
\hline
PedirBus & $grant=1$ & EnviarAddr & $Bus\_req=1$ \\
\hline
EnviarAddr & $DevSel=0$ & ErrorBus & $Frame=1$, $MC\_send\_addr\_ctrl=1$ \\
\hline
EnviarAddr & non-cacheable & TransferPalabra & $Frame=1$, $block\_addr=0$ \\
\hline
EnviarAddr & cacheable y dirty & CopyBack & $Frame=1$, $send\_dirty=1$, $MC\_bus\_Write=1$ \\
\hline
EnviarAddr & cacheable y clean & FetchBloque & $Frame=1$, $MC\_bus\_Read=1$, $block\_addr=1$ \\
\hline
CopyBack & no termina bloque & CopyBack & $MC\_send\_data=1$, $count\_enable=1$ \\
\hline
CopyBack & termina bloque & EnviarAddr & $last\_word=1$, $Block\_copied\_back=1$, $Update\_dirty=1$ \\
\hline
FetchBloque & no termina bloque & FetchBloque & $MC\_WE0/1$, $mux\_origen=1$, $count\_enable=1$ \\
\hline
FetchBloque & termina bloque & FinAcceso & $last\_word=1$, $MC\_tags\_WE=1$, $inc\_m=1$ \\
\hline
TransferPalabra & $TRDY=0$ & TransferPalabra & $Frame=1$ \\
\hline
TransferPalabra & $TRDY=1$ & FinAcceso & lectura: $mux\_output=01$, escritura: $MC\_send\_data=1$ \\
\hline
ErrorBus & siempre & Inicio & $ready=1$, $next\_error\_state=memory\_error$, $load\_addr\_error=1$ \\
\hline
FinAcceso & siempre & Inicio & cierre del acceso \\
\hline
\end{longtable}

\end{document}
